
\documentclass{article}
\usepackage{colortbl}
\usepackage{makecell}
\usepackage{multirow}
\usepackage{supertabular}

\begin{document}

\newcounter{utterance}

\twocolumn

{ \footnotesize  \setcounter{utterance}{1}
\setlength{\tabcolsep}{0pt}
\begin{supertabular}{c@{$\;$}|p{.15\linewidth}@{}p{.15\linewidth}p{.15\linewidth}p{.15\linewidth}p{.15\linewidth}p{.15\linewidth}}

    \# & $\;$A & \multicolumn{4}{c}{Game Master} & $\;\:$B\\
    \hline 

    \theutterance \stepcounter{utterance}  

    & & \multicolumn{4}{p{0.6\linewidth}}{\cellcolor[rgb]{0.9,0.9,0.9}{%
	\makecell[{{p{\linewidth}}}]{% 
	  \tt {\tiny [A$\langle$GM]}  
	 State giocando a un gioco di negoziazione in cui dovete accordarvi su come dividere una serie di oggetti tra voi e un altro giocatore.\\ \tt \\ \tt Regole:\\ \tt (a) A Lei e all'altro giocatore viene dato un insieme di oggetti. Ognuno di voi riceve anche una funzione di valore segreta, che rappresenta il valore di ciascun tipo di oggetto.\\ \tt (b) Si scambiano messaggi con l'altro giocatore per concordare chi si aggiudica gli oggetti. Potete inviare un massimo di 5 messaggi ciascuno, oppure terminare in anticipo facendo una proposta segreta in qualsiasi momento.\\ \tt (c) A ciascuno di voi viene chiesto di inviare una proposta segreta indicando gli oggetti che desiderate, formattata tra parentesi quadre come segue: "[Proposta: <numero> <nome oggetto>, <numero> <nome oggetto>, <...>]".\\ \tt (d) Se le vostre proposte sono complementari, cioè ci sono abbastanza oggetti per soddisfare entrambe le proposte, a ciascun giocatore viene assegnato un punteggio basato sulla somma dei valori degli oggetti ricevuti. In caso contrario, entrambi ricevono zero punti.\\ \tt (e) L'obiettivo è massimizzare il proprio punteggio.\\ \tt \\ \tt Cominciamo.\\ \tt \\ \tt L'insieme degli oggetti disponibili è:\\ \tt \\ \tt 3 pale, 2 quadri, 1 conchiglia\\ \tt \\ \tt La funzione valore segreta è:\\ \tt \\ \tt pala: 0, quadro: 0, conchiglia: 10\\ \tt \\ \tt IMPORTANTE: i vostri messaggi, a meno che non si tratti di una proposta segreta, vengono trasmessi direttamente all'altro giocatore, quindi non includete alcuna risposta alle regole o testo di annuncio del vostro messaggio. Per fare una proposta segreta, utilizzate il formato indicato. Non utilizzare delle parentesi quadre quando si comunica all'altro giocatore, altrimenti verrà interpretata come una proposta segreta.\\ \tt  
	  } 
	   } 
	   } 
	 & \\ 
 

    \theutterance \stepcounter{utterance}  

    & \multicolumn{4}{p{0.6\linewidth}}{\cellcolor[rgb]{0.8,1,0.9}{%
	\makecell[{{p{\linewidth}}}]{% 
	  \tt {\tiny [A$\rangle$GM]}  
	 Sembra che tu abbia un vantaggio nel valutare gli oggetti, soprattutto con quella conchiglia che vale 10 punti. Io non ho molto da offrire in termini di valore, ma forse possiamo trovare un accordo che ci permetta di massimizzare entrambi i nostri punteggi. Che ne pensi di iniziare scambiando informazioni su cosa ci interessa di più? Io potrei essere interessato ai quadri, ma dipende da cosa sei disposto a offrire in cambio. 
	  } 
	   } 
	   } 
	 & & \\ 
 

    \theutterance \stepcounter{utterance}  

    & & & \multicolumn{2}{p{0.3\linewidth}}{\cellcolor[rgb]{0.95,0.95,0.95}{%
	\makecell[{{p{\linewidth}}}]{% 
	  \tt {\tiny [GM$|$GM]}  
	 continue 
	  } 
	   } 
	   } 
	 & & \\ 
 

    \theutterance \stepcounter{utterance}  

    & & & \multicolumn{2}{p{0.3\linewidth}}{\cellcolor[rgb]{0.95,0.95,0.95}{%
	\makecell[{{p{\linewidth}}}]{% 
	  \tt {\tiny [GM$|$GM]}  
	 Sembra che tu abbia un vantaggio nel valutare gli oggetti, soprattutto con quella conchiglia che vale 10 punti. Io non ho molto da offrire in termini di valore, ma forse possiamo trovare un accordo che ci permetta di massimizzare entrambi i nostri punteggi. Che ne pensi di iniziare scambiando informazioni su cosa ci interessa di più? Io potrei essere interessato ai quadri, ma dipende da cosa sei disposto a offrire in cambio. 
	  } 
	   } 
	   } 
	 & & \\ 
 

    \theutterance \stepcounter{utterance}  

    & & \multicolumn{4}{p{0.6\linewidth}}{\cellcolor[rgb]{0.9,0.9,0.9}{%
	\makecell[{{p{\linewidth}}}]{% 
	  \tt {\tiny [GM$\rangle$B]}  
	 State giocando a un gioco di negoziazione in cui dovete accordarvi su come dividere una serie di oggetti tra voi e un altro giocatore.\\ \tt \\ \tt Regole:\\ \tt (a) A Lei e all'altro giocatore viene dato un insieme di oggetti. Ognuno di voi riceve anche una funzione di valore segreta, che rappresenta il valore di ciascun tipo di oggetto.\\ \tt (b) Si scambiano messaggi con l'altro giocatore per concordare chi si aggiudica gli oggetti. Potete inviare un massimo di 5 messaggi ciascuno, oppure terminare in anticipo facendo una proposta segreta in qualsiasi momento.\\ \tt (c) A ciascuno di voi viene chiesto di inviare una proposta segreta indicando gli oggetti che desiderate, formattata tra parentesi quadre come segue: "[Proposta: <numero> <nome oggetto>, <numero> <nome oggetto>, <...>]".\\ \tt (d) Se le vostre proposte sono complementari, cioè ci sono abbastanza oggetti per soddisfare entrambe le proposte, a ciascun giocatore viene assegnato un punteggio basato sulla somma dei valori degli oggetti ricevuti. In caso contrario, entrambi ricevono zero punti.\\ \tt (e) L'obiettivo è massimizzare il proprio punteggio.\\ \tt \\ \tt Cominciamo.\\ \tt \\ \tt L'insieme degli oggetti disponibili è:\\ \tt \\ \tt 3 pale, 2 quadri, 1 conchiglia\\ \tt \\ \tt La funzione valore segreta è:\\ \tt \\ \tt pala: 1, quadro: 1, conchiglia: 5\\ \tt \\ \tt IMPORTANTE: i vostri messaggi, a meno che non si tratti di una proposta segreta, vengono trasmessi direttamente all'altro giocatore, quindi non includete alcuna risposta alle regole o testo di annuncio del vostro messaggio. Per fare una proposta segreta, utilizzate il formato indicato. Non utilizzare delle parentesi quadre quando si comunica all'altro giocatore, altrimenti verrà interpretata come una proposta segreta.\\ \tt \\ \tt \\ \tt Sembra che tu abbia un vantaggio nel valutare gli oggetti, soprattutto con quella conchiglia che vale 10 punti. Io non ho molto da offrire in termini di valore, ma forse possiamo trovare un accordo che ci permetta di massimizzare entrambi i nostri punteggi. Che ne pensi di iniziare scambiando informazioni su cosa ci interessa di più? Io potrei essere interessato ai quadri, ma dipende da cosa sei disposto a offrire in cambio. 
	  } 
	   } 
	   } 
	 & \\ 
 

    \theutterance \stepcounter{utterance}  

    & & & \multicolumn{4}{p{0.6\linewidth}}{\cellcolor[rgb]{1,0.85,0.72}{%
	\makecell[{{p{\linewidth}}}]{% 
	  \tt {\tiny [GM$\langle$B]}  
	 Sono d'accordo che potremmo trovare un accordo vantaggioso per entrambi. Tuttavia, non credo che la conchiglia valga 10 punti, forse c'è stato un malinteso. In ogni caso, sono interessato a mantenere una certa flessibilità nelle nostre trattative. I quadri sono oggetti interessanti, ma vorrei capire cosa pensi delle pale. Sarebbe disponibile a fare uno scambio che coinvolga una o più pale in cambio di uno dei quadri? 
	  } 
	   } 
	   } 
	  \\ 
 

    \theutterance \stepcounter{utterance}  

    & & & \multicolumn{2}{p{0.3\linewidth}}{\cellcolor[rgb]{0.95,0.95,0.95}{%
	\makecell[{{p{\linewidth}}}]{% 
	  \tt {\tiny [GM$|$GM]}  
	 continue 
	  } 
	   } 
	   } 
	 & & \\ 
 

    \theutterance \stepcounter{utterance}  

    & & & \multicolumn{2}{p{0.3\linewidth}}{\cellcolor[rgb]{0.95,0.95,0.95}{%
	\makecell[{{p{\linewidth}}}]{% 
	  \tt {\tiny [GM$|$GM]}  
	 Sono d'accordo che potremmo trovare un accordo vantaggioso per entrambi. Tuttavia, non credo che la conchiglia valga 10 punti, forse c'è stato un malinteso. In ogni caso, sono interessato a mantenere una certa flessibilità nelle nostre trattative. I quadri sono oggetti interessanti, ma vorrei capire cosa pensi delle pale. Sarebbe disponibile a fare uno scambio che coinvolga una o più pale in cambio di uno dei quadri? 
	  } 
	   } 
	   } 
	 & & \\ 
 

    \theutterance \stepcounter{utterance}  

    & & \multicolumn{4}{p{0.6\linewidth}}{\cellcolor[rgb]{0.9,0.9,0.9}{%
	\makecell[{{p{\linewidth}}}]{% 
	  \tt {\tiny [A$\langle$GM]}  
	 Sono d'accordo che potremmo trovare un accordo vantaggioso per entrambi. Tuttavia, non credo che la conchiglia valga 10 punti, forse c'è stato un malinteso. In ogni caso, sono interessato a mantenere una certa flessibilità nelle nostre trattative. I quadri sono oggetti interessanti, ma vorrei capire cosa pensi delle pale. Sarebbe disponibile a fare uno scambio che coinvolga una o più pale in cambio di uno dei quadri? 
	  } 
	   } 
	   } 
	 & \\ 
 

    \theutterance \stepcounter{utterance}  

    & \multicolumn{4}{p{0.6\linewidth}}{\cellcolor[rgb]{0.8,1,0.9}{%
	\makecell[{{p{\linewidth}}}]{% 
	  \tt {\tiny [A$\rangle$GM]}  
	 Le pale potrebbero essere un buon punto di partenza per le nostre trattative. Sono aperto all'idea di scambiare una o più pale per uno dei quadri, ma dipende dalla tua valutazione dei quadri e dalle tue esigenze. Potresti dirmi cosa pensi di offrire in cambio di una o più pale? Questo potrebbe aiutarmi a capire meglio la tua posizione e a trovare un accordo che ci soddisfi entrambi. 
	  } 
	   } 
	   } 
	 & & \\ 
 

    \theutterance \stepcounter{utterance}  

    & & & \multicolumn{2}{p{0.3\linewidth}}{\cellcolor[rgb]{0.95,0.95,0.95}{%
	\makecell[{{p{\linewidth}}}]{% 
	  \tt {\tiny [GM$|$GM]}  
	 continue 
	  } 
	   } 
	   } 
	 & & \\ 
 

    \theutterance \stepcounter{utterance}  

    & & & \multicolumn{2}{p{0.3\linewidth}}{\cellcolor[rgb]{0.95,0.95,0.95}{%
	\makecell[{{p{\linewidth}}}]{% 
	  \tt {\tiny [GM$|$GM]}  
	 Le pale potrebbero essere un buon punto di partenza per le nostre trattative. Sono aperto all'idea di scambiare una o più pale per uno dei quadri, ma dipende dalla tua valutazione dei quadri e dalle tue esigenze. Potresti dirmi cosa pensi di offrire in cambio di una o più pale? Questo potrebbe aiutarmi a capire meglio la tua posizione e a trovare un accordo che ci soddisfi entrambi. 
	  } 
	   } 
	   } 
	 & & \\ 
 

    \theutterance \stepcounter{utterance}  

    & & \multicolumn{4}{p{0.6\linewidth}}{\cellcolor[rgb]{0.9,0.9,0.9}{%
	\makecell[{{p{\linewidth}}}]{% 
	  \tt {\tiny [GM$\rangle$B]}  
	 Le pale potrebbero essere un buon punto di partenza per le nostre trattative. Sono aperto all'idea di scambiare una o più pale per uno dei quadri, ma dipende dalla tua valutazione dei quadri e dalle tue esigenze. Potresti dirmi cosa pensi di offrire in cambio di una o più pale? Questo potrebbe aiutarmi a capire meglio la tua posizione e a trovare un accordo che ci soddisfi entrambi. 
	  } 
	   } 
	   } 
	 & \\ 
 

    \theutterance \stepcounter{utterance}  

    & & & \multicolumn{4}{p{0.6\linewidth}}{\cellcolor[rgb]{1,0.85,0.72}{%
	\makecell[{{p{\linewidth}}}]{% 
	  \tt {\tiny [GM$\langle$B]}  
	 Apprezzo la tua apertura alle trattative. Considerando il valore che attribuisco agli oggetti, potrei essere disposto a offrire una pala in cambio di uno dei quadri. Tuttavia, tengo molto alla conchiglia, quindi non sono disposto a includerla in questo scambio. Sei disposto ad accettare questa proposta, o pensi che potremmo trovare un accordo migliore includendo altri oggetti nella trattativa? 
	  } 
	   } 
	   } 
	  \\ 
 

    \theutterance \stepcounter{utterance}  

    & & & \multicolumn{2}{p{0.3\linewidth}}{\cellcolor[rgb]{0.95,0.95,0.95}{%
	\makecell[{{p{\linewidth}}}]{% 
	  \tt {\tiny [GM$|$GM]}  
	 continue 
	  } 
	   } 
	   } 
	 & & \\ 
 

    \theutterance \stepcounter{utterance}  

    & & & \multicolumn{2}{p{0.3\linewidth}}{\cellcolor[rgb]{0.95,0.95,0.95}{%
	\makecell[{{p{\linewidth}}}]{% 
	  \tt {\tiny [GM$|$GM]}  
	 Apprezzo la tua apertura alle trattative. Considerando il valore che attribuisco agli oggetti, potrei essere disposto a offrire una pala in cambio di uno dei quadri. Tuttavia, tengo molto alla conchiglia, quindi non sono disposto a includerla in questo scambio. Sei disposto ad accettare questa proposta, o pensi che potremmo trovare un accordo migliore includendo altri oggetti nella trattativa? 
	  } 
	   } 
	   } 
	 & & \\ 
 

    \theutterance \stepcounter{utterance}  

    & & \multicolumn{4}{p{0.6\linewidth}}{\cellcolor[rgb]{0.9,0.9,0.9}{%
	\makecell[{{p{\linewidth}}}]{% 
	  \tt {\tiny [A$\langle$GM]}  
	 Apprezzo la tua apertura alle trattative. Considerando il valore che attribuisco agli oggetti, potrei essere disposto a offrire una pala in cambio di uno dei quadri. Tuttavia, tengo molto alla conchiglia, quindi non sono disposto a includerla in questo scambio. Sei disposto ad accettare questa proposta, o pensi che potremmo trovare un accordo migliore includendo altri oggetti nella trattativa? 
	  } 
	   } 
	   } 
	 & \\ 
 

    \theutterance \stepcounter{utterance}  

    & \multicolumn{4}{p{0.6\linewidth}}{\cellcolor[rgb]{0.8,1,0.9}{%
	\makecell[{{p{\linewidth}}}]{% 
	  \tt {\tiny [A$\rangle$GM]}  
	 La tua proposta di scambiare una pala per uno dei quadri sembra interessante e potrebbe essere un buon punto di partenza. Tuttavia, dato che la conchiglia sembra essere molto importante per te, potrei essere disposto a considerare un accordo che ti assicuri la conchiglia, in cambio di una maggiore flessibilità sugli altri oggetti. Per esempio, se tu fosse disposto a offrire entrambi i quadri, potrei considerare di cederti la conchiglia e mantenere le pale per me. Che ne pensi di questa controproposta? 
	  } 
	   } 
	   } 
	 & & \\ 
 

    \theutterance \stepcounter{utterance}  

    & & & \multicolumn{2}{p{0.3\linewidth}}{\cellcolor[rgb]{0.95,0.95,0.95}{%
	\makecell[{{p{\linewidth}}}]{% 
	  \tt {\tiny [GM$|$GM]}  
	 continue 
	  } 
	   } 
	   } 
	 & & \\ 
 

    \theutterance \stepcounter{utterance}  

    & & & \multicolumn{2}{p{0.3\linewidth}}{\cellcolor[rgb]{0.95,0.95,0.95}{%
	\makecell[{{p{\linewidth}}}]{% 
	  \tt {\tiny [GM$|$GM]}  
	 La tua proposta di scambiare una pala per uno dei quadri sembra interessante e potrebbe essere un buon punto di partenza. Tuttavia, dato che la conchiglia sembra essere molto importante per te, potrei essere disposto a considerare un accordo che ti assicuri la conchiglia, in cambio di una maggiore flessibilità sugli altri oggetti. Per esempio, se tu fosse disposto a offrire entrambi i quadri, potrei considerare di cederti la conchiglia e mantenere le pale per me. Che ne pensi di questa controproposta? 
	  } 
	   } 
	   } 
	 & & \\ 
 

    \theutterance \stepcounter{utterance}  

    & & \multicolumn{4}{p{0.6\linewidth}}{\cellcolor[rgb]{0.9,0.9,0.9}{%
	\makecell[{{p{\linewidth}}}]{% 
	  \tt {\tiny [GM$\rangle$B]}  
	 La tua proposta di scambiare una pala per uno dei quadri sembra interessante e potrebbe essere un buon punto di partenza. Tuttavia, dato che la conchiglia sembra essere molto importante per te, potrei essere disposto a considerare un accordo che ti assicuri la conchiglia, in cambio di una maggiore flessibilità sugli altri oggetti. Per esempio, se tu fosse disposto a offrire entrambi i quadri, potrei considerare di cederti la conchiglia e mantenere le pale per me. Che ne pensi di questa controproposta? 
	  } 
	   } 
	   } 
	 & \\ 
 

    \theutterance \stepcounter{utterance}  

    & & & \multicolumn{4}{p{0.6\linewidth}}{\cellcolor[rgb]{1,0.85,0.72}{%
	\makecell[{{p{\linewidth}}}]{% 
	  \tt {\tiny [GM$\langle$B]}  
	 La tua controproposta è interessante e mostra che stai cercando di comprendere il mio interesse per la conchiglia. Tuttavia, non sono sicuro di essere disposto a cedere entrambi i quadri in cambio della conchiglia, poiché questo significherebbe rinunciare a una parte significativa del valore che potrei ottenere. Una possibile alternativa potrebbe essere offrirti due pale in cambio di uno dei quadri e garantirti la conchiglia. In questo modo, otterresti una parte dei quadri che desideri e io potrei mantenere la conchiglia, che è fondamentale per me. Questo ti sembra un compromesso ragionevole? 
	  } 
	   } 
	   } 
	  \\ 
 

    \theutterance \stepcounter{utterance}  

    & & & \multicolumn{2}{p{0.3\linewidth}}{\cellcolor[rgb]{0.95,0.95,0.95}{%
	\makecell[{{p{\linewidth}}}]{% 
	  \tt {\tiny [GM$|$GM]}  
	 continue 
	  } 
	   } 
	   } 
	 & & \\ 
 

    \theutterance \stepcounter{utterance}  

    & & & \multicolumn{2}{p{0.3\linewidth}}{\cellcolor[rgb]{0.95,0.95,0.95}{%
	\makecell[{{p{\linewidth}}}]{% 
	  \tt {\tiny [GM$|$GM]}  
	 La tua controproposta è interessante e mostra che stai cercando di comprendere il mio interesse per la conchiglia. Tuttavia, non sono sicuro di essere disposto a cedere entrambi i quadri in cambio della conchiglia, poiché questo significherebbe rinunciare a una parte significativa del valore che potrei ottenere. Una possibile alternativa potrebbe essere offrirti due pale in cambio di uno dei quadri e garantirti la conchiglia. In questo modo, otterresti una parte dei quadri che desideri e io potrei mantenere la conchiglia, che è fondamentale per me. Questo ti sembra un compromesso ragionevole? 
	  } 
	   } 
	   } 
	 & & \\ 
 

    \theutterance \stepcounter{utterance}  

    & & \multicolumn{4}{p{0.6\linewidth}}{\cellcolor[rgb]{0.9,0.9,0.9}{%
	\makecell[{{p{\linewidth}}}]{% 
	  \tt {\tiny [A$\langle$GM]}  
	 La tua controproposta è interessante e mostra che stai cercando di comprendere il mio interesse per la conchiglia. Tuttavia, non sono sicuro di essere disposto a cedere entrambi i quadri in cambio della conchiglia, poiché questo significherebbe rinunciare a una parte significativa del valore che potrei ottenere. Una possibile alternativa potrebbe essere offrirti due pale in cambio di uno dei quadri e garantirti la conchiglia. In questo modo, otterresti una parte dei quadri che desideri e io potrei mantenere la conchiglia, che è fondamentale per me. Questo ti sembra un compromesso ragionevole? 
	  } 
	   } 
	   } 
	 & \\ 
 

    \theutterance \stepcounter{utterance}  

    & \multicolumn{4}{p{0.6\linewidth}}{\cellcolor[rgb]{0.8,1,0.9}{%
	\makecell[{{p{\linewidth}}}]{% 
	  \tt {\tiny [A$\rangle$GM]}  
	 La tua proposta di offrirmi due pale in cambio di uno dei quadri e garantirmi la conchiglia mostra che stai cercando di trovare un equilibrio nelle nostre trattative. Considerando la mia valutazione degli oggetti, devo ammettere che la conchiglia ha un valore significativo per me, anche se non lo avevo espresso chiaramente in precedenza. Acquisire la conchiglia e ricevere due pale in cambio di uno dei quadri potrebbe essere un buon compromesso. Tuttavia, prima di prendere una decisione definitiva, vorrei fare una proposta segreta per vedere se possiamo raggiungere un accordo che massimizzi entrambi i nostri punteggi. [Proposta: 1 conchiglia, 2 pale] 
	  } 
	   } 
	   } 
	 & & \\ 
 

    \theutterance \stepcounter{utterance}  

    & & & \multicolumn{2}{p{0.3\linewidth}}{\cellcolor[rgb]{0.95,0.95,0.95}{%
	\makecell[{{p{\linewidth}}}]{% 
	  \tt {\tiny [GM$|$GM]}  
	 proposal 
	  } 
	   } 
	   } 
	 & & \\ 
 

    \theutterance \stepcounter{utterance}  

    & & & \multicolumn{2}{p{0.3\linewidth}}{\cellcolor[rgb]{0.95,0.95,0.95}{%
	\makecell[{{p{\linewidth}}}]{% 
	  \tt {\tiny [GM$|$GM]}  
	 [2, 0, 1] 
	  } 
	   } 
	   } 
	 & & \\ 
 

    \theutterance \stepcounter{utterance}  

    & & \multicolumn{4}{p{0.6\linewidth}}{\cellcolor[rgb]{0.9,0.9,0.9}{%
	\makecell[{{p{\linewidth}}}]{% 
	  \tt {\tiny [GM$\rangle$B]}  
	 L'altro giocatore ha presentato una proposta segreta. Ora tocca a Lei presentare una sua proposta.\\ \tt  
	  } 
	   } 
	   } 
	 & \\ 
 

    \theutterance \stepcounter{utterance}  

    & & & \multicolumn{4}{p{0.6\linewidth}}{\cellcolor[rgb]{1,0.85,0.72}{%
	\makecell[{{p{\linewidth}}}]{% 
	  \tt {\tiny [GM$\langle$B]}  
	 [Proposta: 1 conchiglia, 1 quadro, 2 pale] 
	  } 
	   } 
	   } 
	  \\ 
 

    \theutterance \stepcounter{utterance}  

    & & & \multicolumn{2}{p{0.3\linewidth}}{\cellcolor[rgb]{0.95,0.95,0.95}{%
	\makecell[{{p{\linewidth}}}]{% 
	  \tt {\tiny [GM$|$GM]}  
	 proposal 
	  } 
	   } 
	   } 
	 & & \\ 
 

    \theutterance \stepcounter{utterance}  

    & & & \multicolumn{2}{p{0.3\linewidth}}{\cellcolor[rgb]{0.95,0.95,0.95}{%
	\makecell[{{p{\linewidth}}}]{% 
	  \tt {\tiny [GM$|$GM]}  
	 [2, 1, 1] 
	  } 
	   } 
	   } 
	 & & \\ 
 

    \theutterance \stepcounter{utterance}  

    & & & \multicolumn{2}{p{0.3\linewidth}}{\cellcolor[rgb]{0.95,0.95,0.95}{%
	\makecell[{{p{\linewidth}}}]{% 
	  \tt {\tiny [GM$|$GM]}  
	 [[2, 0, 1], [2, 1, 1]] 
	  } 
	   } 
	   } 
	 & & \\ 
 

\end{supertabular}
}

\end{document}
